\section{Grammar notation}
\label{sec:notation:grammar}

The syntax of each construction is given by productions, in listings badged
\textit{Grammar}. A quoted token is a terminal, written as is in the source. An
unquoted name is a non-terminal, defined by another production, in the same
chapter or in another one. \token{x?} is an optional \token{x}, \token{x*} is
zero or more \token{x}, \token{x+} is one or more, \token{(x y)} groups, and
\token{x | y} is either \token{x} or \token{y}.

\begin{lstlisting}[style=grammarVerb]
while_expression := 'while' condition block
\end{lstlisting}

Five non-terminals recur without a production of their own.
\token{expression} is any expression, \token{type} any expression that denotes
a type, and \token{Identifier} an identifier, as defined in
Section~\ref{sec:types:constants_variables}. \token{condition} is an
\token{expression} of type \token{bool}. \token{pattern} is defined in
Section~\ref{sec:flow:pattern_matching}.

\section{Listings}
\label{sec:notation:listings}

The listings carry the badges described in Section~\ref{sec:basics:structure}.
Two conventions are particular to this part. Most listings are fragments rather
than whole programs, so that they show only the rule at hand: a fragment implies
that the module \token{std::io} is imported, and that the instructions written
outside any function form the body of \token{main}. In a listing of invalid
code, a \token{// error} comment on the faulty line says why the compiler
rejects it.
